% Chapter 2 -- The genesis compiler INITECHC.  Sources (Law 1):
% seed/initechc.c (the option loop: --emit-asm, -o, -, unknown options;
% every diagnostic and exit code), seed/README.md, seed/lexer.c (the
% keyword table: `case' is not reserved by the seed, os/tps/PARSER-TRACE.md
% records the divergence), seed/typecheck.c (whole-program collection:
% later routines visible without forward, os/tps/SYMBOLS-DUMP.md s7).  The
% switches are registered (KIND switch) against tps_gate inventory, which
% reads them from initechc.c.  The genesis compiler's agreement with Free
% Pascal is test-seed-fpc-diff; its reproducibility test-seed-repro.
\Gates{test-seed-fpc-diff, test-seed-repro, test-tps-lex-build, test-pubs-tps}
\chapter{The genesis compiler}
\ix{INITECHC}
\ix{genesis compiler}

\file{INITECHC} is a Pascal compiler written in C.  It runs on the
development system, reads one Pascal source file and writes 80386 assembly
language in the same form as Turbo Initech.  It exists to compile the
Turbo Initech source the first time.

\section{The command}
\SyntaxLine{initechc \opt{--emit-asm} \opt{-o \ph{file}} \opt{\ph{source} | -}}

\begin{BorTable}{INITECHC options}{p{26mm}p{72mm}}
\Hd{Option} & \Hd{Effect}\EndHd
\SwitchItem{--emit-asm} & write assembly language; without it,
  \file{INITECHC} checks the program and prints \code{parse OK} and a
  listing of the program's structure\\
\SwitchItem{-o} \textit{file} & write the assembly language to \textit{file}
  instead of the standard output\\
\SwitchItem{-} & read the source from the standard input (also the default
  when no source file is named)\\
\end{BorTable}

\section{Messages and exit codes}
\ix{INITECHC, messages}
An error in the program is reported on the standard error output as

\begin{TypeIn}
source:line:col: error: message
\end{TypeIn}

with \code{<stdin>} as the \textit{source} when the program was read from
the standard input, and \file{INITECHC} ends with exit code 1.  Like Turbo
Initech, it reports only the first error.  A problem with the command or
its files ends it with exit code 2 and one of these messages:
\code{-o needs an argument}, \code{unknown option:} followed by the option
and a usage line, \code{multiple input files}, \code{cannot open}
followed by the file name, \code{read error}, or \code{cannot open output}
followed by the file name.  A failure to close the output file is reported
as \code{error closing} and ends with exit code 1.  A successful run ends
with exit code 0.

\section{The language of the genesis compiler}
\ix{INITECHC, language}
The genesis compiler accepts the language of Turbo Initech, with two
known differences:

\begin{itemize}
\item \code{case} is not a reserved word of the genesis compiler.
\item The genesis compiler reads the whole program before it checks it, so
  a routine may call a routine declared after it without \rw{forward}.
\end{itemize}

Turbo Initech's own source is written so that both compilers accept it.
The genesis compiler's output is checked against the Free Pascal compiler
on a set of programs (\code{make test-seed-fpc-diff}), and every program in
the Turbo Initech manuals is compiled and run by it as well.
